\documentclass[12pt]{article}
\usepackage[utf8]{inputenc}
\usepackage[T1]{fontenc}
\usepackage{hyperref}
\title{Tami Language Resources for K-12 Curriculum Development}
\author{Compiled by Claude Sonnet 5 (claude-sonnet-5) for Andr\'as Kornai}
\date{}
\begin{document}
\maketitle

\section{Introduction}
Tami (ISO 639-3 tmy) is an Austronesian language of Papua New Guinea. It is spoken on the Tami Islands and in a few villages at the tip of the Huon Peninsula, by about 2,100 people (2010). This report rests on a single round of web searches. Little material was found, and nothing suitable for schools.

\section{Primary Educational Resources}
No primers, textbooks or school programmes for Tami turned up in the search results.

\section{Linguistic and Reference Materials}
\begin{itemize}
\item Wikipedia overview of the language (SVO word order, five vowels, labiovelar and prenasalized consonants): \url{https://en.wikipedia.org/wiki/Tami_language}
\item ANU Archives record for a \emph{Vocabulary of Tami language} by Bamler, in German, with a concise grammar in Sio. It is a catalogue record, and I did not confirm that the content is digitized: \url{https://archivescollection.anu.edu.au/index.php/vocabulary-of-tami-language-by-bamler-in-german-and-concise-grammar-in-sio}
\item Numeral data (quinary system): \url{https://lingweb.eva.mpg.de/channumerals/Tami.htm}
\end{itemize}

\section{Religious and Community Resources}
Joshua Project lists a Tami people-group profile that mentions Bible portions. I did not verify which portions exist or where they are held: \url{https://joshuaproject.net/people_groups/15233/PP}

\section{Digital and Audio Resources}
Omniglot has an alphabet chart and a short sample text: \url{https://www.omniglot.com/writing/tami.htm}. No audio or video was found.

\section{Curriculum Development Context}
The search found no modern grammar, no dictionary usable by teachers, and no mother-tongue education programme. A curriculum would have to start with community-led documentation and orthography work. The old Bamler vocabulary could serve as a historical source.
\end{document}